\chapter{Evaluation}\label{chap:evaluation}

This chapter explains the software and heuristics we are evaluating, and the benchmarks and hardware we tested them on. We then present the raw data obtained from these evaluations, along with plots for visualisation, and discussion of inferences we can make from these results.


\section{Evaluated Software}


We integrated our heuristics into the VHPOP POCL planning system by~\cite{Younes2003VHPOPAndAdditiveHeuristic}. The $\mathit{ILP}(\mathcal{P}^+)$ heuristic could not solve any non-trivial problems from our test problems, and so we mention results for the $\mathit{ILP_c}(\mathcal{P}^+)$ heuristic, which we will refer to as the ILP\textsubscript{c} heuristic. We compare the performance of planning using this heuristic against the state of the art heuristic within VHPOP, the ADDR heuristic, which is a version of the additive heuristic by~\cite{Younes2003VHPOPAndAdditiveHeuristic} which accounts for multiple instances of the same action. We also compare results for the state-based planner Fast Downward by~\cite{Helmert2006FastDownward}.\\


For ILP\textsubscript{c}, we use the LCFR-Loc flaw selection strategy included natively within VHPOP~\citep{Younes2003VHPOPAndAdditiveHeuristic}. LCFR-Loc solves causal threats before open conditions, which are preconditions not supported by a causal link. Among open conditions, LCFR-Loc favours local open conditions, which are open conditions for the same step as the most recently supported open condition. To determine the ordering of causal threats, open conditions which apply to the same step, or new open conditions when there are no local open conditions left, LCFR-Loc uses a \textit{Least Cost Flaw Repair} strategy. Least Cost Flaw Repair prioritises flaws with smallest set of possible refinements to solve them, thus minimising the branching factor of the search space. Remaining ties are broken by LIFO, preferring newly discovered flaws, and then by chance if LIFO is also tied (for example when two open preconditions are added at the same time by the same new plan step).\\


We chose the additive heuristic as a comparison as it is currently the best performing POCL heuristic, and is the only state-of-the-art POCL heuristic which is implemented natively within VHPOP. VHPOP implements multiple variations of the additive heuristic, and we chose ADDR as it is the most informed and effective version. For state-based heuristics to compare to, we use the add heuristic by~\cite{Geffner2000Add}, which the additive heuristic for POCL planning is an adaptation of, and the FF heuristic by~\cite{Hoffmann2001FFPlanning}, as it is a delete-relaxation based heuristic.\\


\section{Benchmarks and Setup}


We evaluate all heuristics on the first 10 instances of each STRIPS domain from the 3rd International Planning Competition (IPC3). IPC3 was chosen as it is the planning competition that VHPOP was originally evaluated on, and so VHPOP is built to handle the version of PDDL input files used in this competition. To evaluate the effectiveness of our heuristic for detecting unsolvable problems, we also tested the first 10 instances of STRIPS domains from the 1st Unsolvability International Planning Competition (UIPC1), held in 2016. This is the only well recognised international planning competition for unsolvable problems, as it has not yet been held again, and so is the best option for evaluating our heuristics on unsolvable domains. Our heuristic is built for grounded propositional planning, where all variables are booleans, but the domains in these evaluation sets are not initially grounded. Hence, all problems in our evaluation set are grounded using CPDDL~\citep{Fiser2020FamGroups}, which also infers some mutex propositions and prunes the grounded representation accordingly. For some domains in UIPC1, this grounding process leads to the domains becoming too large to be realistically solvable. We hence evaluate against the 6 domains from UIPC1 with the smallest grounded domain file sizes. Fast Downward accepts these grounded PDDL files as input, but translates them to SAS+ format before solving them. This translation time is included in the evaluation times for Fast Downward heuristics.\\


For each heuristic on each instance, we gave the planning system 10 minutes of CPU time on a 12 core Intel i7-12700H processor (2.30 GHz) with 8GB of RAM. The ILP\textsubscript{c} heuristic is solved using version 22.1.1 of the CPLEX solver, restricted to 1 CPU core. The evaluation was run on a laptop within Windows Subsystem for Linux using the Ubuntu 22 distribution.\\


\section{Empirical Results}


Our results are divided into runtime, nodes generated (i.e. the size of the search space), and coverage (i.e. the number of problems solved within the time limit) on both solvable and unsolvable domains. Since our report is more focused on the potential benefits of more informed POCL heuristics, and our implementation hasn't been fully optimised for runtime, our focus is on the number of nodes generated by our heuristic. We are also interested in determining whether the extra information our heuristic utilises is effective for determining that problems are unsolvable, where other planners might struggle with exponentially growing state trees, and so we will also focus on the coverage results for unsolvable instances.\\

The coverage of solvable domains is shown in Table~\ref{fig:solvableCoverage}. As expected, the state-based heuristics have the highest coverage, and the ADDR heuristic has higher coverage than the ILP\textsubscript{c} heuristic. There were no instances which a state-based heuristic could not solve but ADDR could, and there were no instances ADDR could not solve that ILP\textsubscript{c} could. Table~\ref{fig:unsolvableCoverage} shows the coverage of unsolvable domains. It is expected that the state-based heuristics have higher coverage than the ADDR heuristic, but we had hoped to find that our ILP\textsubscript{c} heuristic has coverage over instances that the other heuristics did not. However, our heuristic performed very poorly on unsolvable instances, even compared to its relative weakness on solvable instances. All instances which ILP\textsubscript{c} determined to be unsolvable were also determined to be unsolvable by all other heuristics in significantly less time. There are likely multiple factors which led to this result. Firstly, the ILP takes significantly longer per search node on unsolvable search nodes compared to solvable search nodes, as the ILP solver needs to perform a significantly larger amount of work to determine that there are no feasible solutions to an ILP. As unsolvable instances typically require the evaluation of more unsolvable search nodes, this disadvantage could outweigh the advantages of being more informed. Additionally, it could be that the additional relaxations in ILP\textsubscript{c} over our standard delete-relaxed heuristic are detrimental to the benefits on unsolvable problems, while the restrictive constraints of our standard delete-relaxed heuristic still make it infeasible for use on any domains due to the evaluation time per search node. A more efficient implementation of the standard delete-relaxation heuristic could potentially still outperform other heuristics on some unsolvable domains, and our results do not conclusively prove that our hypothesis on the benefits of NP-complete heuristics for determining unsolvability is incorrect. More research is still required.


\begin{table}[h]
    \centering
    \begin{tabular}{lccccc}
        \hline
        \textbf{Domain} & $\mathbf{n}$ & \textbf{ILP\textsubscript{c}} & \textbf{ADDR} & \textbf{FD-Add} & \textbf{FD-FF} \\ 
        \hline
        Depots & 10 & 1 & 4 & 8 & 9 \\ 
        DriverLog & 10 & 4 & 10 & 10 & 10 \\ 
        FreeCell & 10 & 0 & 1 & 10 & 10 \\ 
        Rovers & 10 & 4 & 7 & 10 & 10 \\ 
        Satellite & 10 & 3 & 10 & 10 & 10 \\ 
        ZenoTravel & 10 & 2 & 10 & 10 & 10 \\ 
        \hline
        Total & 60 & 14 & 42 & 58 & 59 \\ 
        \hline
    \end{tabular}
    \caption{Coverage table for the first 10 instances of IPC3 domains for VHPOP and Fast Downward heuristics.}
    \label{fig:solvableCoverage}
\end{table}


\begin{table}[h]
    \centering
    \begin{tabular}{lccccc}
        \hline
        \textbf{Domain} & $\mathbf{n}$ & \textbf{ILP\textsubscript{c}} & \textbf{ADDR} & \textbf{FD-Add} & \textbf{FD-FF} \\ 
        \hline
        Bottleneck & 10 & 1 & 10 & 10 & 10 \\ 
        CaveDiving & 10 & 0 & 0 & 6 & 6 \\ 
        ChessboardPebbling & 10 & 0 & 0 & 10 & 6 \\ 
        Pegsol & 10 & 0 & 0 & 10 & 10 \\ 
        PegsolRow5 & 10 & 1 & 3 & 5 & 5 \\ 
        SlidingTiles & 10 & 0 & 0 & 10 & 10 \\ 
        \hline
        Total & 60 & 2 & 13 & 51 & 47 \\ 
        \hline
    \end{tabular}
    \caption{Coverage table for the first 10 instances of 6 UIPC1 domains for VHPOP and Fast Downward heuristics.}
    \label{fig:unsolvableCoverage}
\end{table}


Since ILP\textsubscript{c} could only solve two unsolvable instances, analysis of runtime and generated nodes for unsolvable instances is not relevant, as two instances is too few for any useful conclusions to be drawn apart from the low coverage. Hence, we now focus solely on the IPC 3 domains. Figure~\ref{fig:runtime} shows a log scaled runtime comparison of ILP\textsubscript{c} and ADDR on instances in IPC3 which they both solved before timeout. ADDR has a lower runtime on all instances, which is unsurprising, but the magnitude of this difference is significant. Due to a multiple order of magnitude difference in runtime, it is clear why ILP\textsubscript{c} being more informed does not pay out, as the runtime cost to evaluate the heuristic is far too high. Significant efficiency improvements to our ILP model would be required for it to be a viable alternative to other POCL heuristics. Breakdowns of runtime for each instance are provided in Table~\ref{fig:depotsTime} through to Table~\ref{fig:zenoTravelTime}. These show that while state-based heuristics in Fast Downward typically outperform ADDR in VHPOP, there are some instances where ADDR is fastest; although most of these cases can be attributed to the translation time from PDDL to SAS+ by Fast Downward.


\begin{figure}[!hbt]
	\centering
	\begin{tikzpicture}
	\begin{axis}[
		%title={Number of Generated Search Nodes (Log Scale)},
		xlabel style={anchor=north,yshift=-1em},  % Move x-axis label down
		ylabel style={anchor=east,yshift=1.5em,xshift=4.5em},   % Move y-axis label to the left
		xlabel={ILP\textsubscript{c} (ms)},
		ylabel={ADDR (ms)},
		scale only axis,
		height=10cm,
		width=10cm,
		xmode=log,
		ymode=log,
		xmin=1,
		xmax=1000000,
		ymin=1,
		ymax=1000000
	]
	\addplot+[
		mark=*,
		mark size=3pt,
		color=black,
		only marks]
	table[]
	{data/runtime.dat};
	\addplot[black] coordinates {(1,1) (1000000,1000000)};
	\end{axis}
	\end{tikzpicture}
	\caption{Runtime comparison between ILP\textsubscript{c} and ADDR on IPC3 domains, log scale.}
        \label{fig:runtime}
\end{figure}


\begin{table}[p]
    \centering
    \begin{tabular}{lcccc}
        \hline
        \textbf{Instance} & \textbf{ILP\textsubscript{c}} & \textbf{ADDR} & \textbf{FD-Add} & \textbf{FD-FF} \\ 
        \hline
        Depots1 & 151.37 & 0.01 & 0.65 & 0.64\\ 
        Depots2 & - & 0.3 & 4.94 & 4.81\\ 
        Depots3 & - & - & 4.38 & 4.57\\ 
        Depots4 & - & - & 4.31 & 4.34\\ 
        Depots5 & - & - & 4.59 & 4.45\\ 
        Depots6 & - & - & - & -\\ 
        Depots7 & - & 1.3 & 4.42 & 4.32\\ 
        Depots8 & - & - & 8.07 & 4.57\\ 
        Depots9 & - & - & - & 156.65\\ 
        Depots10 & - & 0.86 & 4.88 & 4.64\\ 
        \hline
    \end{tabular}
    \caption{Solve time in seconds for the first 10 instances of the Depots domain from IPC3. A dash represents a timeout.}
    \label{fig:depotsTime}
\end{table}


\begin{table}[p]
    \centering
    \begin{tabular}{lcccc}
        \hline
        \textbf{Instance} & \textbf{ILP\textsubscript{c}} & \textbf{ADDR} & \textbf{FD-Add} & \textbf{FD-FF} \\ 
        \hline
        DriverLog1 & 14.07 & 0.00 & 0.08 & 0.09\\ 
        DriverLog2 & - & 0.04 & 0.12 & 0.10\\ 
        DriverLog3 & 54.20 & 0.01 & 0.11 & 0.10\\ 
        DriverLog4 & - & 0.17 & 0.13 & 0.13\\ 
        DriverLog5 & - & 0.01 & 0.10 & 0.13\\ 
        DriverLog6 & 142.48 & 0.01 & 0.12 & 0.12\\ 
        DriverLog7 & 90.78 & 0.01 & 0.12 & 0.13\\ 
        DriverLog8 & - & 0.03 & 0.14 & 0.14\\ 
        DriverLog9 & - & 0.01 & 0.18 & 0.16\\ 
        DriverLog10 & - & 0.01 & 0.20 & 0.20\\
        \hline
    \end{tabular}
    \caption{Solve time in seconds for the first 10 instances of the DriverLog domain from IPC3. A dash represents a timeout.}
    \label{fig:driverLogTime}
\end{table}


\begin{table}[p]
    \centering
    \begin{tabular}{lcccc}
        \hline
        \textbf{Instance} & \textbf{ILP\textsubscript{c}} & \textbf{ADDR} & \textbf{FD-Add} & \textbf{FD-FF} \\ 
        \hline
        FreeCell1 & - & 0.03 & 5.26 & 5.27\\ 
        FreeCell2 & - & - & 6.82 & 6.57\\ 
        FreeCell3 & - & - & 7.63 & 7.69 \\ 
        FreeCell4 & - & - & 8.14 & 8.11\\ 
        FreeCell5 & - & - & 10.04 & 9.93  \\ 
        FreeCell6 & - & - & 12.37 & 12.41\\ 
        FreeCell7 & - & - & 12.44 & 12.28  \\ 
        FreeCell8 & - & - & 15.44 & 14.81 \\ 
        FreeCell9 & - & - & 15.5 & 14.73 \\ 
        FreeCell10 & - & - & 18.58 & 18.2\\ 
        \hline
    \end{tabular}
    \caption{Solve time in seconds for the first 10 instances of the FreeCell domain from IPC3. A dash represents a timeout.}
    \label{fig:freeCellTime}
\end{table}


\begin{table}[p]
    \centering
    \begin{tabular}{lcccc}
        \hline
        \textbf{Instance} & \textbf{ILP\textsubscript{c}} & \textbf{ADDR} & \textbf{FD-Add} & \textbf{FD-FF} \\ 
        \hline
        Rovers1 & 19.2 & 0.04 & 0.09 & 0.09\\ 
        Rovers2 & 1.82 & 0 & 0.08 & 0.07\\ 
        Rovers3 & 12.15 & 0 & 0.09 & 0.1\\ 
        Rovers4 & 3.78 & 0 & 0.08 & 0.1\\ 
        Rovers5 & - & - & 0.11 & 0.11\\ 
        Rovers6 & - & - & 0.11 & 0.1\\ 
        Rovers7 & - & 0.04 & 0.1 & 0.11 \\ 
        Rovers8 & - & 1.56 & 0.14 & 0.12\\ 
        Rovers9 & - & 4.12 & 0.13 & 0.11\\ 
        Rovers10 & - & - & 0.15 & 0.15\\ 
        \hline
    \end{tabular}
    \caption{Solve time in seconds for the first 10 instances of the Rovers domain from IPC3. A dash represents a timeout.}
    \label{fig:roversTime}
\end{table}


\begin{table}[p]
    \centering
    \begin{tabular}{lcccc}
        \hline
        \textbf{Instance} & \textbf{ILP\textsubscript{c}} & \textbf{ADDR} & \textbf{FD-Add} & \textbf{FD-FF} \\ 
        \hline
        Satellite1 & 14.22 & 0.01 & 0.09 & 0.09\\ 
        Satellite2 & 181.54 & 0.00 & 0.08 & 0.11\\ 
        Satellite3 & 95.12 & 0.01 & 0.10 & 0.09\\ 
        Satellite4 & - & 9.32 & 0.11 & 0.12\\ 
        Satellite5 & - & 0.01 & 0.12 & 0.12 \\ 
        Satellite6 & - & 0.01 & 0.14 & 0.13\\ 
        Satellite7 & - & 8.29 & 0.16 & 0.15 \\ 
        Satellite8 & - & 0.01 & 0.22 & 0.22\\ 
        Satellite9 & - & 0.01 & 0.24 & 0.25 \\ 
        Satellite10 & - & 0.01 & 0.30 & 0.30\\ 
        \hline
    \end{tabular}
    \caption{Solve time in seconds for the first 10 instances of the Satellite domain from IPC3. A dash represents a timeout.}
    \label{fig:satelliteTime}
\end{table}


\begin{table}[p]
    \centering
    \begin{tabular}{lcccc}
        \hline
        \textbf{Instance} & \textbf{ILP\textsubscript{c}} & \textbf{ADDR} & \textbf{FD-Add} & \textbf{FD-FF} \\ 
        \hline
        ZenoTravel1 & 6.07 & 0.00 & 0.10 & 0.10\\ 
        ZenoTravel2 & - & 0.03 & 0.08 & 0.09\\ 
        ZenoTravel3 & 169.88 & 0.01 & 0.13 & 0.12\\ 
        ZenoTravel4 & - & 0.76 & 0.13 & 0.13\\ 
        ZenoTravel5 & - & 1.20 & 0.15 & 0.17\\ 
        ZenoTravel6 & - & 0.02 & 0.15 & 0.15\\ 
        ZenoTravel7 & - & 0.74 & 0.17 & 0.17\\ 
        ZenoTravel8 & - & 0.97 & 0.30 & 0.28\\ 
        ZenoTravel9 & - & 1.06 & 0.28 & 0.29\\ 
        ZenoTravel10 & - & 5.46 & 0.29 & 0.28\\ 
        \hline
    \end{tabular}
    \caption{Solve time in seconds for the first 10 instances of the ZenoTravel domain from IPC3. A dash represents a timeout.}
    \label{fig:zenoTravelTime}
\end{table}


Figure~\ref{fig:nodes} shows a log scaled comparison of the number of nodes generated by ILP\textsubscript{c} and ADDR on instances in IPC3 which they both solved before timeout. It is clear that ILP\textsubscript{c} is a more informed heuristic, as it generated less nodes on all instances. However, the difference in nodes generated was not enough to overcome the runtime weaknesses of our heuristic. Despite this, a clear and significant difference in generated nodes is an important and promising result. Without counting relaxations this difference would likely be more pronounced, and so a more efficiently modelled delete-relaxation heuristic could provide enough of an information benefit that heuristic runtime could be offset by a smaller search tree. Breakdowns of generated nodes for each instance are provided in Table~\ref{fig:depotsNodes} through to Table~\ref{fig:zenoTravelNodes}. These tables show that state-based heuristics usually generated less search nodes than POCL heuristics. In most circumstances this would not be the case, but VHPOP is operating on propositional PDDL and fast downward translates that propositional PDDL into SAS+ before solving. SAS+ allows for multi-values variables and other more expressive techniques which give it an advantage over propositional PDDL. Hence, the result that state-based methods generate less nodes is not evidence that state-based planning creates smaller search trees, as the opposite is true, but instead that SAS+ generates smaller search trees than propositional PDDL.


\begin{figure}[!hbt]
\centering
\begin{tikzpicture}
\begin{axis}[
	%title={Number of Generated Search Nodes (Log Scale)},
	xlabel style={anchor=north,yshift=-1em},  % Move x-axis label down
	ylabel style={anchor=east,yshift=1.5em,xshift=4.5em},   % Move y-axis label to the left
	xlabel={ILP\textsubscript{c}},
	ylabel={ADDR},
	scale only axis,
	height=10cm,
	width=10cm,
	xmode=log,
    ymode=log,
    xmin=1,
    xmax=1000000,
    ymin=1,
    ymax=1000000
]
\addplot+[
	mark=*,
	mark size=3pt,
	color=black,
	only marks]
table[]
{data/nodes.dat};
\addplot[black] coordinates {(1,1) (1000000,1000000)};
\end{axis}
\end{tikzpicture}
\caption{Generated nodes comparison between ILP\textsubscript{c} and ADDR on IPC3 domains, log scale.}
\label{fig:nodes}
\end{figure}


\begin{table}[p]
    \centering
    \begin{tabular}{lrrrr}
        \hline
        \textbf{Instance} & \textbf{ILP\textsubscript{c}} & \textbf{ADDR} & \textbf{FD-Add} & \textbf{FD-FF} \\ 
        \hline
        Depots1 & 651 & 4721 & 83 & 136 \\
        Depots2 & - & 289521 & 195 & 382 \\
        Depots3 & - & - & 796 & 13438 \\
        Depots4 & - & - & 3975 & 84177 \\
        Depots5 & - & - & 17443 & 54639 \\
        Depots6 & - & - & - & - \\ 
        Depots7 & - & 1125861 & 3788 & 14310 \\ 
        Depots8 & - & - & 700926 & 46652 \\
        Depots9 & - & - & - & 28081239 \\ 
        Depots10 & - & 901451 & 89359 & 51639 \\
        \hline
    \end{tabular}
    \caption{Number of nodes generated for the first 10 instances of the Depots domain from IPC3. A dash represents a timeout.}
    \label{fig:depotsNodes}
\end{table}


\begin{table}[p]
    \centering
    \begin{tabular}{lrrrr}
        \hline
        \textbf{Instance} & \textbf{ILP\textsubscript{c}} & \textbf{ADDR} & \textbf{FD-Add} & \textbf{FD-FF} \\ 
        \hline
        DriverLog1 & 271 & 771 & 61 & 61 \\
        DriverLog2 & - & 53231 & 32945 & 23803 \\ 
        DriverLog3 & 1141 & 20311 & 178 & 185 \\ 
        DriverLog4 & - & 300611 & 43202 & 32478 \\ 
        DriverLog5 & - & 6181 & 565 & 33976 \\
        DriverLog6 & 2011 & 6261 & 6404 & 7360 \\ 
        DriverLog7 & 1501 & 3611 & 375 & 486 \\
        DriverLog8 & - & 26341 & 14914 & 71203 \\ 
        DriverLog9 & - & 4171 & 18164 & 3780 \\
        DriverLog10 & - & 21711 & 856 & 564 \\
        \hline
    \end{tabular}
    \caption{Number of nodes generated for the first 10 instances of the DriverLog domain from IPC3. A dash represents a timeout.}
    \label{fig:driverLogNodes}
\end{table}


\begin{table}[p]
    \centering
    \begin{tabular}{lrrrr}
        \hline
        \textbf{Instance} & \textbf{ILP\textsubscript{c}} & \textbf{ADDR} & \textbf{FD-Add} & \textbf{FD-FF} \\ 
        \hline
        FreeCell1 & - & 28901 & 52 & 57 \\
        FreeCell2 & - & - & 136 & 100 \\ 
        FreeCell3 & - & - & 145 & 148 \\
        FreeCell4 & - & - & 1130 & 350 \\
        FreeCell5 & - & - & 1434 & 337 \\ 
        FreeCell6 & - & - & 2102 & 979 \\ 
        FreeCell7 & - & - & 1228 & 5272 \\
        FreeCell8 & - & - & 3244 & 1380 \\ 
        FreeCell9 & - & - & 3547 & 1308 \\ 
        FreeCell10 & - & - & 4473 & 1691 \\ 
        \hline
    \end{tabular}
    \caption{Number of nodes generated for the first 10 instances of the FreeCell domain from IPC3. A dash represents a timeout.}
    \label{fig:freeCellNodes}
\end{table}


\begin{table}[p]
    \centering
    \begin{tabular}{lrrrr}
        \hline
        \textbf{Instance} & \textbf{ILP\textsubscript{c}} & \textbf{ADDR} & \textbf{FD-Add} & \textbf{FD-FF} \\ 
        \hline
        Rovers1 & 2341 & 150871 & 98 & 98 \\
        Rovers2 & 381 & 1681 & 68 & 68 \\ 
        Rovers3 & 1371 & 1711 & 189 & 136 \\ 
        Rovers4 & 531 & 931 & 67 & 73 \\
        Rovers5 & - & - & 778 & 777 \\ 
        Rovers6 & - & - & 2386 & 2446 \\ 
        Rovers7 & - & 90231 & 587 & 598 \\ 
        Rovers8 & - & 2861541 & 1666 & 1035 \\ 
        Rovers9 & - & 6524031 & 1764 & 1550 \\ 
        Rovers10 & - & - & 3624 & 3208 \\
        \hline
    \end{tabular}
    \caption{Number of nodes generated for the first 10 instances of the Rovers domain from IPC3. A dash represents a timeout.}
    \label{fig:roversNodes}
\end{table}


\begin{table}[p]
    \centering
    \begin{tabular}{lrrrr}
        \hline
        \textbf{Instance} & \textbf{ILP\textsubscript{c}} & \textbf{ADDR} & \textbf{FD-Add} & \textbf{FD-FF} \\ 
        \hline
        Satellite1 & 1681 & 1431 & 70 & 84 \\ 
        Satellite2 & 6861 & 2881 & 384 & 156 \\ 
        Satellite3 & 4451 & 1821 & 208 & 208 \\
        Satellite4 & - & 22055171 & 562 & 622 \\ 
        Satellite5 & - & 5471 & 1187 & 1393 \\ 
        Satellite6 & - & 6071 & 819 & 1152 \\ 
        Satellite7 & - & 14454561 & 3663 & 2183 \\
        Satellite8 & - & 11461 & 5186 & 3351 \\ 
        Satellite9 & - & 13891 & 3764 & 5414 \\
        Satellite10 & - & 15371 & 5256 & 7643 \\
        \hline
    \end{tabular}
    \caption{Number of nodes generated for the first 10 instances of the Satellite domain from IPC3. A dash represents a timeout.}
    \label{fig:satelliteNodes}
\end{table}


\begin{table}[p]
    \centering
    \begin{tabular}{lrrrr}
        \hline
        \textbf{Instance} & \textbf{ILP\textsubscript{c}} & \textbf{ADDR} & \textbf{FD-Add} & \textbf{FD-FF} \\ 
        \hline
        ZenoTravel1 & 141 & 731 & 7 & 7 \\
        ZenoTravel2 & - & 130831 & 35 & 46 \\
        ZenoTravel3 & 3851 & 9431 & 174 & 184 \\ 
        ZenoTravel4 & - & 2457891 & 128 & 264 \\
        ZenoTravel5 & - & 3348611 & 255 & 305 \\ 
        ZenoTravel6 & - & 38241 & 485 & 411 \\
        ZenoTravel7 & - & 1870181 & 249 & 348 \\ 
        ZenoTravel8 & - & 2151051 & 430 & 1389 \\ 
        ZenoTravel9 & - & 1366201 & 1041 & 1357 \\ 
        ZenoTravel10 & - & 6992631 & 1912 & 2687 \\
        \hline
    \end{tabular}
    \caption{Number of nodes generated for the first 10 instances of the ZenoTravel domain from IPC3. A dash represents a timeout.}
    \label{fig:zenoTravelNodes}
\end{table}


Overall, these results indicate that the use of ILP-based heuristics for finding unsolvable problems is not promising, and a modified approach would be required to better exploit potential advantages of ILP-based heuristics for unsolvable domains. However, there are no conclusive results that rule out our claims about the potential benefits of NP-hard POCL heuristics on unsolvable domains. Additionally, it is clear that our ILP\textsubscript{c} heuristic is more informed than the ADDR heuristic, and this advantage in reducing the size of the search tree provides evidence that more runtime efficient variations of our heuristics could be effective.